% Edit the following

% Personal information
\newcommand{\myname}{Oscar Amo Grau} % Your name (e.g. "John Smith")
\newcommand{\mytitle}{Robotics Software Engineer} % Your title (e.g. "Applicant")
\newcommand{\myemail}{oscaramog@gmail.com} % Your email (e.g. "me@mydomain.com")
\newcommand{\mylinkedin}{oscaramograu} % Your LinkedIn profile name (e.g. "myname")
\newcommand{\myphone}{+49 15757815031} % Your phone number (e.g. "123.456.7890")
\newcommand{\mylocation}{Munich, Germany} % Your general location (e.g. "PNW, USA")
\newcommand{\recipient}{Hiring Manager} % The recipient (e.g. "Hiring Manager")
\newcommand{\greeting}{Dear} % The greeting to use (e.g. "Dear")
\newcommand{\closer}{Kind Regards} % The closer to use (e.g. "Kind Regards")

% Company information
% Left blank on purpose: the base letter names no employer. These are filled in
% per application, in the copy under applications/.
\newcommand{\company}{} % The company (e.g. "Microsoft Corporation")
\newcommand{\city}{} % The company's city (e.g. "Redmond")
\newcommand{\team}{} % The team, if the posting names one (e.g. "Perception")
\newcommand{\position}{} % The role as the posting names it (e.g. "Robotics Software Engineer")

% Why this company, specifically — one paragraph, written fresh for each
% application and left empty here. It renders between the closing paragraphs of
% the letter; nothing appears while it is empty.
%
% Point at something real in the posting or the company: the product, the problem,
% the stack, how they work. Then connect it to something already true in profile/.
% A sentence that could be sent to any other company is worth nothing — cut it.
%
% e.g. \renewcommand{\whyfit}{Your postings mention running perception on the robot
% rather than in a datacentre. That constraint is most of what I have been doing
% for the past two years, and it is the part of the problem I find interesting.}
\newcommand{\whyfit}{}